\par\noindent\textbf{(1)} 波动方程μw\_tt=T∇²w，设w=R(r)e\textasciicircum{}\{imθ\}e\textasciicircum{}\{−iωt\}，角向周期要求m整数。径向方程R''+(1/r)R'+[k²−m²/r²]R=0，其中k²=μω²/T。中心有限性去掉奇异的Y\_m分支，保留J\_m(kr)。
\par\noindent\textbf{(2)} 固定边缘R(a)=0给J\_m(ka)=0；令j\_\{m,n\}为第n个正零点，则ω\_mn=j\_\{m,n\}√(T/μ)/a。对m≠0，cos(mθ)与sin(mθ)二重简并；m=0基模不具角向二重简并。
\par\noindent\textbf{(3)} 径向算符−(1/r)(rR')'+m²R/r²在权重r下自伴。取不同本征解u\_n,u\_p，交叉乘方程相减并积分，端点边界项消失，于是(k\_n²−k\_p²)∫₀\textasciicircum{}a r u\_nu\_pdr=0，n≠p时径向正交，且积分必须带Jacobian权重r。
\par\noindent\textbf{(4)} 最低m=0,n=1零点j\_\{0,1\}≈2.4048256，因此基频ω01=j\_\{0,1\}√(T/μ)/a。增大膜半径使频率按a\textasciicircum{}\{-1\}降低，增大张力按√T提高，增面密度按μ\textasciicircum{}\{-1/2\}降低；边界若非圆形，纯Bessel零点关系失效。
